\documentclass[11pt]{article}
\usepackage[margin=1in]{geometry}
\usepackage{booktabs}
\usepackage{graphicx}
\usepackage{hyperref}
\usepackage{amsmath}

\title{Causal Circuits for Tool Selection in Instruction-Tuned Language Models}
\author{Anonymous Authors}
\date{}

\begin{document}
\maketitle

\begin{abstract}
Large language model agents must decide whether to answer directly or invoke
external tools. We introduce ToolUseCircuitBench, a controlled benchmark for
web search, calculation, Python execution, and direct answering, and study the
internal mechanisms that mediate these choices. Our analysis combines
behavioral evaluation, held-out linear probing, activation geometry,
controlled activation patching, steering, and component ablation. Results and
cross-model replication will be inserted only after all preregistered controls
are complete.
\end{abstract}

\section{Introduction}
Tool-selection errors are consequential but mechanistically underexplored.
This work asks where tool decisions are represented, whether tool types occupy
distinct directions, and which components causally mediate those decisions.

\section{ToolUseCircuitBench}
Describe dataset construction, template-family separation, OOD prompts,
instruction-conflict prompts, label balance, and contamination controls.

\section{Methods}
\subsection{Behavioral evaluation}
\subsection{Residual and MLP probing}
\subsection{Tool directions}
\subsection{Controlled activation patching}
\subsection{Steering}
\subsection{Neuron and attention-head ablation}
\subsection{Statistical analysis}

\section{Results}
\subsection{Behavioral competence and errors}
\subsection{Layer-wise emergence of tool information}
\begin{figure}[t]
  \centering
  \includegraphics[width=.8\linewidth]{figures/layerwise_probe.pdf}
  \caption{Held-out OOD probe performance by layer.}
\end{figure}
\subsection{Tool-direction geometry}
\subsection{Causal patching with matched controls}
\subsection{Necessity and sufficiency}
\subsection{Cross-model replication}

\section{Limitations}
The controlled single-label setting simplifies real multi-step agents.
Quantized inference may alter fine-grained causal effects. Component
selectivity alone is not interpreted as causal evidence.

\section{Conclusion}

\bibliographystyle{plain}
\bibliography{references}
\end{document}
